We give an explicit finite counterexample and then prove the two support assertions.

Let \(\mathcal V=\{A,Y,Z\}\), let \(\mathcal R=\varnothing\), let \(\mathcal X_A=\mathcal X_Y=\{0,1\}\), and let \(\mathcal X_Z=\{0,1,2\}\). Thus every substantive variable is always observed and \(\mathcal O=(A,Y,Z)\). Let \(\mathcal G\) have the directed edge \(A\to Y\), the bidirected edge \(A\leftrightarrow Y\), the isolated vertex \(Z\), and the empty label set. Its directed part is acyclic, as required by \(\mathrm{ass{:}acyclic}\). Because \(\mathcal R=\varnothing\), there is no missingness context in which a label could occur; consequently the empty label set is regular and maximal, as required by \(\mathrm{ass{:}regular}\) and \(\mathrm{ass{:}maximal}\). Assumption \(\mathrm{ass{:}proxy\text{-}consistency}\) is vacuous for this graph.

Let \(\mathcal F\) be the \(C\)-forest induced by \(\{A,Y\}\), and let \(\mathcal F'\) be its subforest induced by \(\{Y\}\). The forest \(\mathcal F\) is bidirected-connected, \(A\) has the unique directed child \(Y\), and \(Y\) has no directed child, so its root set is \(\{Y\}\); \(\mathcal F'\) has the same root set. Moreover,
\[
 \mathcal F\cap\{A\}=\{A\},
 \qquad
 \mathcal F'\cap\{A\}=\varnothing.
\]
The common root \(Y\) is an ancestor of \(Y\) after intervention on \(A\), since a vertex is its own ancestor. Take the empty color assignment and the trace family consisting only of the empty trace. This trace family is prefix closed, and the licensing condition is vacuous because no color occurs. Hence, by \(\mathrm{def{:}colored\text{-}hedge}\),
\[
 H=(\mathcal F,\mathcal F',\varnothing,\{\varnothing\})
 \in \operatorname{CH}(\mathcal G,A,Y).
\]

We next construct two compatible models. Let \(U\sim\operatorname{Bernoulli}(1/2)\) be the exogenous factor shared by \(A\) and \(Y\). Independently, let \(E_A\sim\operatorname{Bernoulli}(1/4)\), let \(E_Y\) be uniform on \(\{1,\ldots,16\}\), and let \(E_Z\) be uniform on \(\{0,1,2\}\). Put
\[
 A=U\mathbin{\oplus}E_A,
 \qquad
 Z=E_Z,
 \qquad
 t_0=\frac14,
 \qquad
 t_1=-\frac14,
\]
where \(\oplus\) denotes addition modulo two and the Bernoulli parameter is the probability of value one. For \(i\in\{0,1\}\), define
\[
 k_i(u,a)
 =16\left[\frac12+t_i\left\{\mathbf 1(u=a)-\frac34\right\}\right],
\]
and in model \(\mathcal M_i\) set
\[
 Y=\mathbf 1\{E_Y\leq k_i(U,A)\}.
\]
These are deterministic finite response tables because the thresholds are integers: \(k_0(u,a)\) is \(9\) when \(u=a\) and \(5\) otherwise, whereas \(k_1(u,a)\) is \(7\) when \(u=a\) and \(11\) otherwise. For each model, the equation for \(A\) depends on the shared factor \(U\), while the equation for \(Y\) depends on both \(A\) and \(U\), because changing the truth value of \(u=a\) changes the threshold. Thus the directed edge \(A\to Y\) and the individual bidirected-edge factor for \(A\leftrightarrow Y\) are active, and no additional edge is introduced. The factors \(U,E_A,E_Y,E_Z\) are mutually independent, so \(\mathrm{ass{:}semimarkovian}\) holds. The empty label set imposes no response-table restriction.

For positivity, define
\[
 \pi_i(u,a)=\frac{k_i(u,a)}{16}.
\]
The displayed threshold values give \(0<\pi_i(u,a)<1\) for every \(i,u,a\). Moreover, \(P(E_A=e)>0\) for both \(e\in\{0,1\}\). Therefore, for every \((a,y,z)\in\mathcal X_A\times\mathcal X_Y\times\mathcal X_Z\), independence and the structural equations give
\[
 \begin{aligned}
 P_{\mathcal M_i}(A=a,Y=y,Z=z)
 &=\frac13\sum_{u=0}^1
 \frac12 P(E_A=a\mathbin{\oplus}u)
 \pi_i(u,a)^y\{1-\pi_i(u,a)\}^{1-y}\\
 &>0.
 \end{aligned}
\]
This verifies \(\mathrm{ass{:}strict\text{-}positivity}\) cell by cell. All alphabets are finite, as required by \(\mathrm{ass{:}finite\text{-}alphabets}\). The graph conditions established above and \(\mathrm{def{:}positive\text{-}rm\text{-}class}\) therefore imply
\[
 \mathcal M_0,\mathcal M_1
 \in\mathfrak M_{+}^{\mathrm{rm}}(\mathcal G).
\]

We now compute the observed laws. For either \(a\in\{0,1\}\),
\[
 P_{\mathcal M_i}(A=a)
 =\sum_{u=0}^1\frac12P(E_A=a\mathbin{\oplus}u)
 =\frac12\left(\frac34+\frac14\right)
 =\frac12,
\]
and
\[
 P_{\mathcal M_i}(U=a,A=a)
 =P(U=a,E_A=0)
 =\frac12\cdot\frac34
 =\frac38.
\]
Division is legitimate because \(P_{\mathcal M_i}(A=a)=1/2>0\), and hence
\[
 P_{\mathcal M_i}(U=a\mid A=a)
 =\frac{3/8}{1/2}
 =\frac34.
\]
Using the response equation for \(Y\) and the conditional law of total probability,
\[
 \begin{aligned}
 P_{\mathcal M_i}(Y=1\mid A=a)
 &=\sum_{u=0}^1
 \left[\frac12+t_i\left\{\mathbf 1(u=a)-\frac34\right\}\right]
 P_{\mathcal M_i}(U=u\mid A=a)\\
 &=\frac12+t_i\left\{P_{\mathcal M_i}(U=a\mid A=a)-\frac34\right\}\\
 &=\frac12.
 \end{aligned}
\]
Since \(Z=E_Z\) is independent and uniform, both models consequently induce
\[
 P_{\mathcal M_i}(A=a,Y=y,Z=z)
 =\frac12\cdot\frac12\cdot\frac13
 =\frac1{12}
\]
for every observed cell. Let \(P_{\mathrm{obs}}\) denote this common uniform, and hence strictly positive, observed law.

Under \(\operatorname{do}(A=a)\), the structural intervention in \(\mathrm{def{:}causal\text{-}query}\) replaces the equation for \(A\) by the constant \(a\) and leaves the exogenous distribution, in particular \(U\sim\operatorname{Bernoulli}(1/2)\), unchanged. Therefore
\[
 \begin{aligned}
 Q_a^{\mathcal M_i}(1)
 &=\sum_{u=0}^1
 \left[\frac12+t_i\left\{\mathbf 1(u=a)-\frac34\right\}\right]P(U=u)\\
 &=\frac12+t_i\left(\frac12-\frac34\right)\\
 &=\frac12-\frac{t_i}{4}.
 \end{aligned}
\]
It follows that
\[
 Q_a^{\mathcal M_0}(1)=\frac7{16},
 \qquad
 Q_a^{\mathcal M_1}(1)=\frac9{16}.
\]
By \(\mathrm{def{:}observational\text{-}fiber}\), both models lie in \(\mathfrak F(\mathcal G,P_{\mathrm{obs}})\), but their interventional laws differ. Thus the displayed colored hedge has two strictly positive observationally equivalent witnesses with different \(Q_a\).

It remains to prove that no universal binary witness can match this observed law and to establish the general support claim. Let \(W\) be always observed and define
\[
 S_W=\{w:P_{\mathrm{obs}}(W=w)>0\}.
\]
In any model inducing \(P_{\mathrm{obs}}\), the observed value of \(W\) is its substantive value, so
\[
 S_W\subseteq\mathcal X_W
 \quad\Longrightarrow\quad
 |\mathcal X_W|\geq |S_W|.
\]
For a partially observed substantive coordinate \(W\), \(\mathrm{ass{:}proxy\text{-}consistency}\) gives \(W^{\dagger}=W\) on \(R_W=1\). The required conditioning event has positive probability under \(\mathrm{ass{:}strict\text{-}positivity}\): by \(\mathrm{ass{:}finite\text{-}alphabets}\), summing the finitely many positive full-data cells with \(R_W=1\) gives \(P_{\mathcal M}(R_W=1)>0\). Consequently
\[
 S_W^{\dagger}
 =\operatorname{supp}_{P_{\mathrm{obs}}}(W^{\dagger}\mid R_W=1)
 \subseteq\mathcal X_W,
\]
and hence
\[
 |\mathcal X_W|\geq |S_W^{\dagger}|.
\]

Under \(\mathrm{ass{:}strict\text{-}positivity}\), both inclusions are equalities at the level of supports. Indeed, for every \(w\in\mathcal X_W\), summing the positive full-data cells whose substantive coordinate is \(W=w\) yields
\[
 P_{\mathcal M}(W=w)>0
\]
when \(W\) is always observed. If \(W\) is partially observed, summing the positive full-data cells with \(W=w\) and \(R_W=1\) gives
\[
 P_{\mathcal M}(W=w,R_W=1)>0.
\]
By \(\mathrm{ass{:}proxy\text{-}consistency}\), this equals
\[
 P_{\mathrm{obs}}(W^{\dagger}=w,R_W=1)>0.
\]
Thus every declared substantive value occurs in the conditional proxy support. Conversely, proxy consistency permits no value outside \(\mathcal X_W\) on \(R_W=1\). We have therefore proved
\[
 |S_W|=|\mathcal X_W|
 \quad\text{for always-observed }W,
 \qquad
 |S_W^{\dagger}|=|\mathcal X_W|
 \quad\text{for partially observed }W.
\]

In the constructed common law, \(S_Z=\{0,1,2\}\). Hence every model matching \(P_{\mathrm{obs}}\) must satisfy \(|\mathcal X_Z|\geq3\), so no such model can make every substantive variable binary. This disproves the universal binary-witness assertion.

Finally, for an arbitrary integer \(m\geq2\), replace \(\mathcal X_Z\) by \(\{0,\ldots,m-1\}\) and take \(E_Z\) uniform on it, leaving the equations for \(A\) and \(Y\) unchanged. The same computation gives two strictly positive models with a common observed law and distinct query values. The forests \(\mathcal F'\subseteq\mathcal F\), their districts and roots, the empty coloring, and the trace family are unchanged because \(Z\) is isolated and belongs to neither forest. On the other hand, the observed support of \(Z\) has cardinality \(m\), so the support argument forces \(|\mathcal X_Z|\geq m\); the displayed construction attains equality. Thus the smallest required \(Z\)-alphabet coordinate is \(m\), while the colored hedge data are fixed. Since \(m\) is arbitrary, the smallest substantive alphabet vector is not determined by the color and district structure of a colored hedge alone. This proves every assertion of the theorem.